% Generated draft: requires mathematical and source review.
\documentclass[11pt]{article}
\usepackage[margin=25mm]{geometry}
\usepackage{fontspec}
\setmainfont{Latin Modern Roman}
\usepackage{amsmath,amssymb,mathtools,unicode-math}
\setmathfont{Latin Modern Math}
\usepackage{enumitem,needspace,xurl,hyperref}
\setlength{\parindent}{0pt}
\setlength{\parskip}{5pt}
\setlength{\emergencystretch}{3em}
\newcommand{\N}{\mathbb{N}}
\newcommand{\Z}{\mathbb{Z}}
\newcommand{\Q}{\mathbb{Q}}
\newcommand{\R}{\mathbb{R}}
\newcommand{\C}{\mathbb{C}}
\newcommand{\F}{\mathbb{F}}
\newcommand{\A}{\mathbb{A}}
\newcommand{\PP}{\mathbb{P}}
\newcommand{\E}{\mathbb{E}}
\newcommand{\eps}{\varepsilon}
\newcommand{\dd}{\,\mathrm{d}}
\begin{document}
\section*{UnsolvedMath ID: 20003334}
\textbf{Source problem code:} catalogue entry
\subsection{Definitions and mathematical statement}
This source record is a setup paragraph, not an autonomous Conjecture 38. The aid says its opening words finish a backward reference and that it ends before the displayed Conjecture 42. The latter statement, the meaning of ``generic,'' and its Gaussian normalization are not supplied and are not invented here.

Let \(X\) be a compact negatively curved Riemannian manifold, \(S^*X\) its unit cotangent bundle, \(g^t\) the geodesic flow, \(\mathcal X\) its generator, and \(m_L\) normalized Liouville probability. Let \(\phi_j\) be \(L^2\)-normalized Laplace eigenfunctions, \(-\Delta\phi_j=\lambda_j\phi_j\), and put \(P=\sqrt{-\Delta}\). For smooth symbols use classical order-zero pseudodifferential quantizations. A symbol pulled back from \(X\) admits multiplication as a chosen quantization. In the surface specialization with area \(4\pi\), Weyl's law is \(N(\Lambda)=\#\{j:\lambda_j\leq\Lambda\}\sim\Lambda\); this surface normalization is not a formula for every dimension.

For \(A,B\in\Psi^0(X)\), with principal symbols \(a,b\) restricted to \(S^*X\), define
\[
 A'=A+i[P,B].
\]
In the standard commutator convention its principal symbol is \(a+\mathcal Xb\), and the aid's exact quantum identity is
\[
 \langle A'\phi_j,\phi_j\rangle=\langle A\phi_j,\phi_j\rangle
\]
for every eigenfunction. For real symbols, self-adjoint choices give the intended real observables. The Liouville mean is also unchanged since \(\int\mathcal Xb\,\dd m_L=0\).

If quantizations of \(a\) and \(a+\mathcal Xb\) are chosen independently, rather than by this exact commutator rule, their difference from the displayed matched construction is of order \(-1\). The aid records the diagonal error
\[
 O(\lambda_j^{-1/2}),
\]
and hence \(O(\lambda_j^{-1/4})\) after multiplication by the proposed surface CLT scale \(\lambda_j^{1/4}\). Constants depend on the fixed operators and geometry. The statement requires the usual classical pseudodifferential calculus; it is not an assertion for arbitrary unrelated quantization procedures.

For the independent classical comparison, center \(a\) by \(a^0=a-\int a\,\dd m_L\) and define
\[
 V_T(a)=\frac1T\left\|\int_0^T a^0\circ g^t\,\dd t\right\|_{L^2(m_L)}^2.
\]
The flow-coboundary identity \(\int_0^T\mathcal Xb\circ g^t\,\dd t=b\circ g^T-b\) yields the reported bound
\[
 \left|\sqrt{V_T(a+\mathcal Xb)}-\sqrt{V_T(a)}\right|
 \leq\frac{2\|b\|_{L^2(m_L)}}{\sqrt T}.
\]
Thus whenever the limiting variance exists and is finite, the two limiting variances agree. When represented by a convergent Green--Kubo correlation integral, that variance is likewise unchanged. Existence or convergence hypotheses for such an integral are not silently omitted.

In particular, a pure coboundary has limiting variance zero. A normalized Gaussian formulation dividing by the square root of that variance must exclude or separately handle this class; it cannot automatically assert a nondegenerate Gaussian there. The commutator and variance identities do not prove a quantum central limit theorem, characterize generic negative curvature, or recover the missing Conjecture 42. They are the supplied invariance lemma and normalization audit within the excerpt's setup.
\subsection{Short English statement}
Adding the specified quantum commutator changes the principal symbol by a flow coboundary while leaving every diagonal eigenfunction expectation unchanged. The classical finite-time variance has an explicit matching bound. Quantization errors vanish at the proposed surface scaling, but no CLT is proved and zero-variance classes need separate treatment.
\subsection{Sources}
\begin{itemize}
\item Local UnsolvedMath registry, permanent ID \(20003334\), exact source code \texttt{catalogue entry}; canonical \texttt{aim-other-notes.json}, \texttt{notes[140]}, AIM domain Other, registry category Miscellaneous, tag \texttt{conjecture}.
\item AIM workshop ``Emerging applications of measure rigidity'': \url{https://aimath.org/WWN/measrigid/measrigid.pdf}. The aid reports that this paragraph bridges a backward reference to Conjecture 38 and the subsequent Conjecture 42; neither full statement is included. The source's exact label describes the assigned fragment. No PDF was fetched.
\item Source-local labels [55] and [44] appear in the backward reference. Full bibliographic entries and authoritative citations for the stated calculus/variance lemma are not supplied. The original literature field is empty.
\item Research attempt 1: result \texttt{context\_only}, problem status at run \texttt{not\_a\_problem}, rigor label \texttt{proved}, novelty confidence \texttt{low}; classification \texttt{OPEN-TRIAGE}, registry status \texttt{open}, suggested difficulty L2. Created and updated 14 August 2026. These are supplied assessments, not acceptance of a missing CLT conjecture.
\item Attribution: Ulam AI / UnsolvedMath Contributors, CC BY 4.0. The machine-generated research aid is not peer reviewed and requires independent verification before citation.
\end{itemize}
\end{document}
